% Part: proof-theory
% Chapter: sequent-calculus
% Section: admissible-derivable

\documentclass[../../../include/open-logic-section]{subfiles}

\begin{document}

\olfileid[pt]{pt}{seq}{adm}

\olsection{Regras admissíveis e deriváveis}

Muitos resultados da teoria da prova dizem respeito a---ou usam---resultados
sobre a questão de tudo o que pode ser !!{prove}d com uma regra específica
também poder ser !!{prove}d sem essa regra. O teorema de eliminação de cortes
é o exemplo mais conhecido. Ele mostra que todo sequente que pode ser !!{prove}d
usando as regras de um sistema de sequentes com~\Cut{} também pode ser !!{prove}d
sem~\Cut. Como esse tipo de propriedade é muito importante, ele tem um
nome.

\begin{defn}
Uma regra~$R$ é \emph{admissível} em um sistema de sequentes se todo sequente
que puder ser demonstrado nesse sistema juntamente com a regra~$R$ também
puder ser demonstrado sem~$R$.
\end{defn}

\begin{ex}\ollabel{ex:land-deriv}
A regra $\LeftR{\land}$ apresenta formas diferentes em \Log{G1c} e
\Log{G3c}. Em \Log{G1c}, há duas versões da regra: uma com o
conjuntivo esquerdo $!A$ da !!{formula} principal~$!A \land !B$ na
premissa e outra com o conjuntivo direito~$!B$:
\[
\Axiom$ !A, \Gamma \fCenter \Delta$
\RightLabel{$\LeftR{\land}_1$}
\UnaryInf$ !A \land !B, \Gamma \fCenter \Delta$
\DisplayProof
\qquad
\Axiom$!B, \Gamma \fCenter \Delta$
\RightLabel{$\LeftR{\land}_2$}
\UnaryInf$!A \land !B, \Gamma \fCenter \Delta$
\DisplayProof
\]
Em contraste, \Log{G3c} tem apenas uma regra:
\[
\Axiom$ !A, !B, \Gamma \fCenter \Delta$
\RightLabel{$\LeftR{\land}_3$}
\UnaryInf$ !A \land !B, \Gamma \fCenter \Delta$
\DisplayProof
\]
Poderíamos agora perguntar: ``A regra de \Log{G3c} $\LeftR{\land}_3$ é
admissível em \Log{G1c}?'' É: podemos simular diretamente toda inferência
que use $\LeftR{\land}_3$ usando somente as regras de~\Log{G1c}. Além de
$\LeftR{\land}_1$ e $\LeftR{\land}_2$, precisaremos também da regra
estrutural~$\LeftR{\Contraction}$:
\[
\Axiom$!A, !B, \Gamma \fCenter \Delta$
\RightLabel{$\LeftR{\land}_1$}
\UnaryInf$!A \land !B, !B, \Gamma \fCenter \Delta$
\RightLabel{$\LeftR{\land}_2$}
\UnaryInf$!A \land !B, !A \land !B, \Gamma \fCenter \Delta$
\RightLabel{\LeftR{\Contraction}}
\UnaryInf$!A \land !B \Gamma \fCenter \Delta$
\DisplayProof
\]
\end{ex}

Simular uma inferências que usa uma regra por meio de outras regras é uma
maneira de mostrar que uma regra é admissível. Mas não é a única maneira,
e algumas regras são admissíveis, porém não podem ser simuladas assim. As
regras que podem são chamadas \emph{deriváveis}.

\begin{defn}
Uma regra~$R$,
\[
\AxiomC{$S_1$}
\AxiomC{$\dots$}
\AxiomC{$S_n$}
\RightLabel{$R$}
\TrinaryInfC{$S$}
\DisplayProof
\]
é chamada \emph{derivável} em um sistema se houver uma !!{proof}
esquemática do sequente~$S$ que use as regras e os axiomas do sistema
juntamente com sequentes iniciais adicionais das formas de~$S_1$, \dots,~$S_n$.
\end{defn}

A !!{proof} apresentada em \olref{ex:land-deriv} mostra que
$\LeftR{\land}_3$ é uma regra derivável em~\Log{G1c}, pois essa !!{proof}
é uma !!{proof} esquemática da conclusão de $\LeftR{\land}_3$ a partir
das premissas de $\LeftR{\land}_3$ como sequentes iniciais, usando apenas
regras de \Log{G1c}. (Ela não usou nenhum dos axiomas de \Log{G1c}.)

\begin{prop}\ollabel{prop:lor-G3-adm}
As regras $\LeftR{\land}$ e $\RightR{\lor}$ de~\Log{G3c} são
deriváveis em~\Log{G1c}.
\end{prop}

\begin{proof}
  A demonstração para $\LeftR{\land}$ é dada em \olref{ex:land-deriv}. A
  demonstração para $\RightR{\lor}$ fica como exercício.
\end{proof}

\begin{prob}
Mostre que a regra $\RightR{\lor}$ de \Log{G3c} é derivável
em~\Log{G1c}.
\end{prob}

\begin{prob}
Suponha que $\liff$ seja um !!{operator} definido: $!A \liff !B$ é uma
abreviação de $(!A \lif !B) \land (!B \lif !A)$. Mostre que as regras
\begin{enumerate}
\item \Axiom$\Gamma, !A, !B \fCenter \Delta$
\Axiom$\Gamma \fCenter \Delta, !A, !B$
\RightLabel{\LeftR{\liff}}
\BinaryInf$\Gamma \fCenter \Delta, !A \liff !B$
\DisplayProof
\item
\Axiom$!A, \Gamma \fCenter \Delta, !B$
\Axiom$!B, \Gamma \fCenter \Delta, !A$
\RightLabel{\RightR{\liff}}
\BinaryInf$\Gamma \fCenter \Delta, !A \liff !B$
\DisplayProof
\end{enumerate}
são deriváveis em (a)~\Log{G1c} e (b)~\Log{G3c}.
\end{prob}


\begin{ex}\ollabel{ex:weak-adm}
A regra $\RightR{\Weakening}$ de \Log{G1c},
\[
\Axiom$ \Gamma \fCenter \Delta$
\RightLabel{\RightR{\Weakening}}
\UnaryInf$ \Gamma \fCenter \Delta, !A$
\DisplayProof
\]
é admissível
em~\Log{G3c}. A ideia da demonstração é simples: suponha que haja
!!a{proof}~$\pi$ da premissa $\Gamma \Sequent \Delta$ em~\Log{G3c}.
Acrescente a !!{formula}~$!A$ ao sucedente de cada sequente
em~$\pi$. Axiomas continuam sendo axiomas; inferências corretas continuam
sendo inferências corretas. A !!{formula} final da nova !!{proof} é $\Gamma \Sequent
\Delta, !A$.

A regra $\RightR{\Weakening}$, contudo, não é derivável
em~\Log{G3c}. Pois suponha que fosse, isto é, suponha que se pudesse encontrar
!!a{proof} de $\Gamma \Sequent
\Delta, !A$ usando somente os axiomas e as regras de~\Log{G3c}, além,
possivelmente, do sequente inicial adicional~$\Gamma \Sequent \Delta$. Se essa
!!{proof} fosse puramente esquemática, como seria necessário para que
$\RightR{\Weakening}$ fosse derivável, poderíamos transformá-la em !!a{proof}
de~$\Sequent !A$ eliminando $\Gamma$ e~$\Delta$. Isso transformaria
os sequentes iniciais adicionais~$\Gamma \Sequent \Delta$ no sequente
vazio~$\quad \Sequent \quad$. Como o sequente vazio não contém
quaisquer !!{formula}s, mas todas as regras de~\Log{G3c} exigem pelo menos
uma !!{formula} ativa em suas premissas, o sequente vazio não pode ser uma
premissa de nenhuma inferência nessa !!{proof} esquemática. Consequentemente, a
!!{proof} só pode ser esquemática se, de fato, não usar o sequente adicional
$\Gamma \Sequent \Delta$ como sequente inicial. Em outras palavras,
ela seria uma derivação esquemática de $\Gamma \Sequent
\Delta, !A$ usando somente os axiomas e as regras de~\Log{G3c}, e teríamos
mostrado que todo sequente dessa forma tem !!a{proof} em~\Log{G3c}.
Mas isso é impossível, pois \Log{G3c} é correto e, portanto, só pode !!{prove}
sequentes que sejam verdadeiros em toda !!{structure}. Tomando $\Gamma = \Delta
= \emptyset$ e $!A \ident \lfalse$, \Log{G3c} teria de !!{prove},
por exemplo, o sequente $\quad \Sequent \lfalse$, o que não ocorre.
\end{ex}

Apresentemos uma boa demonstração indutiva do resultado de que o
enfraquecimento é uma regra admissível em~\Log{G3c}. Na verdade, vale um
resultado ligeiramente mais forte, como mostra o argumento intuitivo: ao
acrescentar $!A$ ao lado direito de cada sequente de~$\pi$, não alteramos a
estrutura da !!{proof}; em particular, não aumentamos a altura. Como isso é
importante, merece uma definição própria.

\begin{defn}
Uma regra~$R$,
\[
\AxiomC{$S_1$}
\AxiomC{$\dots$}
\AxiomC{$S_n$}
\RightLabel{$R$}
\TrinaryInfC{$S$}
\DisplayProof
\]
é \emph{admissível com preservação de !!{height}} (ou
\emph{!!{hp}-admissível}) em um sistema se, sempre que $\Proves[h_i] S_i$
para $i = 1$, \dots,~$n$, então $\Proves_m S$ com $m = \max(h_1, \dots,
h_n)$.
\end{defn}

Observe que, para uma regra ser admissivel com preservação de
!!{height}, basta que, sempre que as premissas~$S_i$ tiverem !!{proof}s~$\pi_i$
com !!{height}~$h_i = \pheight{\pi_i}$, haja !!a{proof}~$\pi$ da
conclusão cuja !!{height}~$\pheight{\pi}$ não seja maior que o
máximo dos~$h_i$. Em particular, se houver somente uma
premissa~$S_1$, basta que $\pheight{\pi} \le \pheight{\pi_1}$.
No caso de $\RightR{\Weakening}$ e $\LeftR{\Weakening}$, temos, de fato,
$\pheight{\pi} = \pheight{\pi_1}$, mas isso não é
necessário.

\begin{prop}\ollabel{prop:weak-G3c-adm}
As regras $\RightR{\Weakening}$ e $\LeftR{\Weakening}$ de \Log{G1c} são
admissíveis com preservação de !!{height} em~\Log{G3c}.
\end{prop}

\begin{proof}
  Temos de mostrar que, se $\Log{G3c} \Proves \Gamma \Sequent \Delta$
  com uma !!{proof}~$\pi$ de !!{height}~$\pheight{\pi} = n$, então há
  uma \Log{G3c}-!!{proof} $\pi'$ de $\Gamma \Sequent \Delta, !A$ com
  $\pheight{\pi'} = n$. Demonstramos isso por indução em~$n$. Se $n=0$,
  então $\pi$ é apenas o axioma $\Gamma \Sequent \Delta$ (isto é, alguma
  $!C$ atômica ocorre tanto em $\Gamma$ quanto em~$\Delta$), e o mesmo vale
  para $\Gamma \Sequent \Delta, !A$. Se $n = \pheight{\pi} >0$, a
  !!{proof}~$\pi$ tem uma última inferência. Distinguimos casos segundo a
  regra utilizada. Damos apenas um exemplo: suponha que a última inferência
  tenha sido $\RightR{\land}$. Então $\Delta = \Delta', !B \land !C$, e as
  sub-!!{proof}s imediatas $\pi_1$ e~$\pi_2$ da última inferência
  têm sequentes finais $\Gamma \Sequent \Delta', !B$ e $\Gamma \Sequent
  \Delta', !C$, respectivamente. Em outras palavras, $\pi$ tem a seguinte forma:
  \[
  \AxiomC{}
  \RightLabel{$\pi_1$}
  \Deduce$\Gamma \fCenter \Delta', !B$
  \AxiomC{}
  \RightLabel{$\pi_2$}
  \Deduce$\Gamma \fCenter \Delta', !C$
  \RightLabel{\RightR{\land}}
  \BinaryInf$\Gamma \fCenter \Delta', !B \land !C$
  \DisplayProof
  \]
  Também temos que $n =
  \max(\pheight{\pi_1}, \pheight{\pi_2}) + 1$. Logo, $\pheight{\pi_1} <
  n$ e $\pheight{\pi_2} < n$, e a hipótese de indução se aplica:
  há !!{proof}s $\pi_1'$ e~$\pi_2'$ de $\Gamma \Sequent
  \Delta', !B, !A$ e $\Gamma \Sequent \Delta', !C, !A$. Usando a regra
  $\RightR{\land}$, obtemos !!a{proof}~$\pi'$ de $\Gamma
  \Sequent \Delta', !B \land !C, !A$:
  \[
  \AxiomC{}
  \RightLabel{$\pi_1'$}
  \Deduce$\Gamma \fCenter \Delta', !B, !A$
  \AxiomC{}
  \RightLabel{$\pi_2'$}
  \Deduce$\Gamma \fCenter \Delta', !C, !A$
  \RightLabel{\RightR{\land}}
  \BinaryInf$\Gamma \fCenter \Delta', !B \land !C, !A$
  \DisplayProof
  \]
  Temos $\pheight{\pi'} =
  \max(\pheight{\pi_1'}, \pheight{\pi_2'}) + 1 = \max(\pheight{\pi_1},
  \pheight{\pi_2}) + 1 = \pheight{\pi}$.
\end{proof}

\olref{prop:weak-G3c-adm} é extremamente útil. Introduziremos uma
notação para sua aplicação repetida: se $\pi$ for !!a{proof}
de~$\Gamma \Sequent \Delta$, então $\pi[\Pi \Sequent \Lambda]$ será o
resultado da aplicação da admissibilidade do enfraquecimento à esquerda com
todas as !!{formula}s em~$\Pi$ e da admissibilidade do enfraquecimento à direita
com todas as !!{formula}s em~$\Lambda$. Trata-se de !!a{proof} de $\Gamma, \Pi
\Sequent \Delta, \Lambda$.

\end{document}
